% Minimal retained preamble for lane-p2 cards taken from
% sources/logic-fl02-conservation-of-ramsey-s-theorem-f/pinned-original-source.tex
% (Conservation of Ramsey's theorem for pairs, arXiv:2402.11616v3).
%
% Copied verbatim from the author's own file:
%   line 2      : \documentclass{article}
%   lines 16-25 : the author's theorem declarations
%   lines 46-88 : the author's macro definitions for the reverse-mathematics principles
%   line 91     : \card
%   lines 104-107: the natural sums and products, which rely on mathabx (source line 103)
%   line 116    : \qt
%
% Omitted author preamble lines and why (nothing was renamed or redefined):
%   line 3 (graphicx), 7 (url), 9 (cancel): no figure, url or cancel use in the retained spans;
%   line 6 (cleveref): the assembler adds its own cleveref machinery, a second load clashes;
%   line 8 (\usepackage[utf8x]{inputenc}): inputenc with the utf8x option fights the forced
%       xelatex engine;
%   lines 28-37 (the rep@theorem machinery and \newreptheorem{maintheorem}): the retained spans
%       use no rep-theorem environment;
%   lines 40-43 (xcolor and the three author comment macros): not used by the retained spans;
%   lines 93-99 (the LS2 font encoding block for \smashtimes based on stix): the host TeX Live has
%       no stix directory and the author himself never uses the result -- line 105 maps \ntimes to
%       \dottimes, and line 110 shows \smashtimes already commented out in the source.
% mathabx is retained because \dotplus, \dottimes and \dotdiv come from it.
\documentclass{article}
\usepackage{amsmath,amssymb,amsthm}
\usepackage{mathabx}
\usepackage{cancel}
\usepackage{cleveref}

\newtheorem{theorem}{Theorem}
\numberwithin{theorem}{section}
\newtheorem{maintheorem}[theorem]{Main Theorem}
\newtheorem{lemma}[theorem]{Lemma}
\newtheorem{question}[theorem]{Question}
\newtheorem{proposition}[theorem]{Proposition}
\newtheorem{statement}[theorem]{Statement}
\newtheorem{remark}[theorem]{Remark}
\newtheorem{definition}[theorem]{Definition}
\newtheorem{corollary}[theorem]{Corollary}

\newcommand{\RCA}{\mathsf{RCA}}
\newcommand{\WKL}{\mathsf{WKL}}
\newcommand{\SRT}{\mathsf{SRT}}
\newcommand{\EM}{\mathsf{EM}}
\newcommand{\SEM}{\mathsf{SEM}}
\newcommand{\ADS}{\mathsf{ADS}}
\newcommand{\SCAC}{\mathsf{SCAC}}
\newcommand{\SADS}{\mathsf{SADS}}
\newcommand{\COH}{\mathsf{COH}}
\newcommand{\ACA}{\mathsf{ACA}}
\newcommand{\qvdash}{\operatorname{{?}{\vdash}}}
\newcommand{\nqvdash}{\operatorname{{?}{\nvdash}}}
\newcommand{\BSig}{\mathsf{B}\Sigma^0}
\newcommand{\ISig}{\mathsf{I}\Sigma^0}
\newcommand{\LSig}{\mathsf{L}\Sigma^0}
\newcommand{\BPi}{\mathsf{B}\Pi^0}
\newcommand{\IPi}{\mathsf{I}\Pi^0}
\newcommand{\LPi}{\mathsf{L}\Pi^0}
\newcommand{\IDel}{\mathsf{I}\Delta^0}
\newcommand{\BCDe}{\mathsf{BC}\Delta^0}
\newcommand{\WF}{\mathsf{WF}}
\newcommand{\RT}{\mathsf{RT}}
\newcommand{\BME}{\mathsf{BME}}
\newcommand{\TT}{\mathsf{TT}}
\renewcommand{\O}{\mathcal{O}}
\renewcommand{\L}{\mathcal{L}}
\newcommand{\C}{\mathcal{C}}
\newcommand{\R}{\mathcal{R}}
\renewcommand{\S}{\mathcal{S}}
\newcommand{\T}{\mathcal{T}}
\renewcommand{\P}{\mathcal{P}}
\newcommand{\M}{\mathcal{M}}
\newcommand{\Nc}{\mathcal{N}}
\renewcommand{\P}{\mathcal{P}}
\newcommand{\Q}{\mathcal{Q}}
\newcommand{\NN}{\mathbb{N}}
\newcommand{\PP}{\mathbb{P}}
\newcommand{\uh}[0]{{\upharpoonright}}
\newcommand{\card}{\operatorname{card}}
\newcommand{\nplus}{\dotplus}
\newcommand{\ntimes}{\dottimes}
\newcommand{\nminus}{\dotdiv}
\newcommand{\nsum}{\dot{\sum}}
\def\qt#1{``#1''}
\title{Every tree condition decides an initial segment of the jump by a finite extension}
\author{Quentin Le Houérou, Ludovic Levy Patey, Keita Yokoyama}
\date{arxiv:2402.11616v3; CC BY 4.0}
\begin{document}
\maketitle
\noindent\textit{Editorial note.} The counted unit is the article's own Lemma lem:wkl-wf-forcing-jump at source lines 666-668 together with the proof environment that begins on the next line (670-680); the pairing is the ordinary adjacent one.\par\medskip
\noindent\textit{Editorial note.} The retained context at source line 645 announces the technical lemma on extendible nodes; the retained context at source lines 647-650 is that lemma, attributed by the author to Fernandes et al, which the proof applies at source line 673; the retained context at source line 663 is the author's explanation of the forcing relation that decides an initial segment of the jump.\par\medskip
\noindent\textit{Editorial note.} The retained context at source lines 640-643, which lists the requirements R\_k, S\_k and T\_k of this forcing, is not retained because the counted lemma does not use them; the tree condition and the extension relation are described in the standard prerequisites instead. The lemma lem:picking-extendible-node-in-tree retained here carries the author's citation of Fernandes et al at source line 647; that citation is part of the article's own statement of the technical lemma and its entry is not retained, so the material carries no bibliography.\par\medskip
\noindent\textit{Editorial note.} The retained preamble reproduces verbatim the author's document class at source line 2, his theorem declarations at source lines 16-25, his macro definitions at source lines 46-88 and 91, his natural-operation macros at source lines 104-107 and his quotation macro at source line 116, and it retains mathabx from source line 103. The LS2 font-encoding block at source lines 93-99 is not retained because the host TeX Live has no stix directory and the author never uses that symbol. The cleveref load at source line 6 is not retained because the assembler adds its own cleveref machinery, which is what the retained \textbackslash{}Cref calls need, and the inputenc load with the utf8x option at source line 8 is not retained because it fights the forced xelatex engine.\par\medskip
\noindent\textit{Editorial note.} The article's conservation theorems, its COH lemmas, its Shore-blocking lemmas and its SEM section lemmas are not counted here and are not used as premises.\par\medskip
\noindent\textit{Editorial note.} The retained context at source lines 647-650 carries the author citation of Fernandes et al inside the title of the technical lemma, and the retained proof uses the author Cref calls. The cited entry is therefore taken verbatim from the author own bibliography file biblio.bib as shipped in the pinned e-print of arXiv:2402.11616v3, attached to this task unedited and complete, and the card is assembled with the bibtex backend. The preamble additionally loads cleveref, which the assembler only adds by itself when it generates navigation stubs.\par\medskip
\noindent\textit{Allowed standard prerequisites (editorial).}\begin{itemize}
\item Tree conditions of the forcing used for WKL. In this section a condition is a pair (sigma, T) consisting of a finite binary string and an M-infinite tree, and an extension (tau, T\_2) of (sigma, T\_1) requires tau to extend sigma and the tree T\_2 to be contained in T\_1; the requirements to be satisfied are listed by the author at source lines 640-643.. Fixes the objects of the counted lemma, namely the condition (sigma, T\_1), its extension and the compatible generics G and Y.
\item Forcing a decided initial segment of the jump. A condition (tau, U) forces the k-th initial segment of the jump of G plus Y to equal a string sigma' when for every index e below the length of sigma' it forces the corresponding value to converge if sigma'(e) is one and to diverge if sigma'(e) is zero; the article records this at source line 663 and notes that the relation is Pi\textasciicircum{}0\_1(Y).. It is the conclusion of the counted lemma, and the proof transfers the two cases of the definition to every generic compatible with the constructed extension.
\item Definability classes and M-finiteness in the model. A set defined by a Pi\textasciicircum{}0\_1(Y) formula is M-finite in the models considered, and Pi\textasciicircum{}0\_1 sets of finite binary strings of a fixed length are M-finite, which is what allows a lexicographically minimal element to be selected; the article uses this throughout this section.. The proof forms the set W of candidate strings, argues that it is Pi\textasciicircum{}0\_1(Y) and M-finite, and then selects its lexicographically minimal element.
\end{itemize}\medskip
We first prove a simple technical lemma saying that there exists arbitrarily large extendible nodes. This lemma even holds over~$\RCA_0^*$ and can be found in various places in the literature.

\begin{lemma}[Fernandes et al~\cite{fernandes2017analysis}]\label[lemma]{lem:picking-extendible-node-in-tree}
Let $\M = (M,S) \models \RCA_0^*$ be a model and $T \subseteq 2^{<M}$ be an $M$-infinite tree in~$\M$.
Then for every~$\ell \in M$, there is some~$\sigma \in T$ of length~$\ell$ such that $\{ \tau \in T : \tau \mbox{ is comparable with } \sigma \}$ is $M$-infinite.
\end{lemma}

In what follows, $(\tau, U) \Vdash (G \oplus Y)' \uh_k = \sigma'$ means that for every~$e < |\sigma'|$, if $\sigma'(e) = 1$ then $(\tau, U) \Vdash \Phi^{G \oplus Y}_e(e)\downarrow$ and if $\sigma'(e) = 0$, then $(\tau, U) \Vdash \Phi^{G \oplus Y}_e(e)\uparrow$. Thus, this relation is $\Pi^0_1(Y)$.

\begin{lemma}\label[lemma]{lem:wkl-wf-forcing-jump}
Let $(\sigma, T_1)$ a condition. For every $k \in M$, there exists some $M$-finite $\sigma'$ and an extension $(\tau, T_2) \leq (\sigma, T_1)$ forcing $(G \oplus Y)' \uh_k = \sigma'$.
\end{lemma}

\begin{proof}
Let $W$ be the set of all $\sigma' \in 2^k$ such that the tree $\{\rho \in T_1 | \sigma'(e) = 0 \to \Phi_e^{\rho \oplus Y}(e) \uparrow\}$ is $M$-infinite. $W$ is $\Pi_0^1(Y)$ and hence $M$-finite, and it is non-empty as it contains the string $1111\dots$.

We can consider $\sigma' \in W$ its lexicographically minimal element, and consider the corresponding tree $T_3$. For every $e < k$ such that $\sigma'(e) = 1$, the minimality of $\sigma'$ gives us that the set of $\tau \in T_3$ such that $\Phi_e^{\tau \oplus Y}(e) \uparrow$ is $M$-finite, so we have a level $h_e$ such that for every $\tau \in T_3 \cap 2^{h_e}$, $\Phi_e^{\tau \oplus Y}(e) \downarrow$. The set $\{e<k|\sigma'(e) = 1\}$ is $M$-finite, so we can consider an upperbound $h$ of all the $h_e$. Finally, by \Cref{lem:picking-extendible-node-in-tree}, there is a $\tau \in T_3 \cap 2^{h}$ such that $T_2 := \{\rho \in T_3 : \rho \textit{ is comparable with } \tau\}$ is $M$-infinite (since $T_3$ is $M$-infinite).

We claim that $(\tau, T_2)$ forces $(G \oplus Y)' \uh_k = \sigma'$. Fix some $e < k$.

Case (a): $\sigma'(e) = 0$. Then $\Phi_e^{\tau \oplus Y}(e) \downarrow$ for every $\tau \in T_2$ (since $T_2 \subseteq T_3$), hence for every $G$ compatible with $(\tau,T_2)$, so $e \notin (G \oplus Y)'$.

Case (b): $\sigma'(e) = 1$. The definition of $\tau$ ensure that $\Phi_e^{\tau \oplus Y}(e) \downarrow$, so $e \in (G \oplus Y)'$. 
\end{proof}

\bibliographystyle{plain}
\bibliography{retained-biblio}
\end{document}
